\subsection{\TextToI Generation}

Diffusion models have revolutionized the field of text-to-image generation. While a comprehensive review of all text-to-image models is beyond the scope of this paper, we will focus on the most relevant ones that have been published, productionized or open-sourced, and thus widely used by a large user base.

The seminal work of latent diffusion models~\citep{rombach2021highresolution} proposes compressing the original image space to latent space using a variational autoencoder, which improves training and inference efficiency, thus popularizing latent-space-based diffusion models. Dalle-3~\citep{dalle3} proposes using GPT to rewrite image captions, reducing noise in curated internet-scale text-image pairs for more effective training. Emu~\citep{dai2023emu} proposes using a higher latent dimension and fine-tuning a pre-trained model with a small, high-quality dataset to exclusively generate high-quality, professional-looking images. Stable Diffusion 3~\citep{sd3} proposes using rectified flow transformers with a multimodal diffusion backbone to improve generation quality.

In MovieGen, we also use a 16-channel variational autoencoder and flow transformers with prompt rewrite to achieve both high visual quality and text alignment.

\subsection{\TextToV Generation}

The swift progress in text-to-image generation has led to substantial improvements in temporally coherent high quality video generation. After the success of diffusion models for image generation~\citep{dhariwal2021diffusion, ramesh2022hierarchical}, they have been vastly used to improve video synthesis~\citep{NEURIPS2022_39235c56}.
Several works introduce zero-shot video generation by enriching the pre-trained text-to-image generation models with motion dynamics~\citep{text2video-zero, wu2023tune}. DirecT2V~\citep{hong2023direct2v} leverages instruction-tuned large language models for zero-shot video creation by dividing user inputs into separate prompts for each frame.


Several other papers propose a cascaded or factorized approach for text-to-video generation.  Imagen-Video~\citep{ho2022imagen} and Make-A-Video~\citep{singer2023makeavideo} trained a deep cascade of spatial and temporal layers via pixel diffusion modeling while many other works focus more on applying diffusion to the latent space of an auto-encoder for more efficiency~\citep{Blattmann2023AlignYL, an2023latentshift, wang2023lavie, wang2023videofactory, modelscope}.
AnimateDiff~\citep{guo2023animatediff} introduces a pre-trained motion module into a pre-trained T2I model. Emu-Video~\citep{emuvideo2023}, Stable Video Diffusion~\citep{sdvideo}, I2VGen-XL~\citep{i2vgen}, Dynamicrafter~\citep{xing2023dynamicrafter}, VideoGen~\citep{li2023videogen}, and VideoCrafter1~\citep{chen2023videocrafter1} add an image as an extra conditioning to the T2V model. Lumiere~\citep{BarTal2024LumiereAS} uses a Space-Time U-Net to generate the full temporal duration of the video at once. SEINE~\citep{chen2023seine} facilitates the smooth integration of shots from diverse scenes and generates videos of various lengths through auto-regressive prediction.

A few papers have studied the role of noise scheduling for more coherent~\citep{ge2023preserve, qiu2023freenoise, luo2023videofusion} and longer~\citep{kim2024fifo} video generation.
While most of the above text-to-video generation models use a U-Net based architecture, Snap-Video~\citep{snapvideo} and \Sora~\citep{sora} show the scalability and out-performance of transformer architectures for diffusion-based video generation.
Latte~\citep{ma2024latte} also uses a DiT instead of the U-Net backbone for text-to-video generation. On the other hand, a couple of works have been focused on transformer models within an auto-regressive framework~\citep{yan2021videogpt, kondratyuk2023videopoet, hong2022cogvideo, wu2022nuwa, villegas2023phenaki, ge2022long}. RIVER~\citep{davtyan2023efficient} uses flow matching for efficient video prediction by conditioning on a small set of past frames in the latent space of a pre-trained VQGAN. In this work, we leverage a Llama3 transformer architecture~\citep{llama3} and train a text-to-video generation model within a flow matching framework~\citep{flow-matching}.


\textbf{Encoding videos into a latent space.} Since their inception~\citep{rombach2021highresolution,Esser2021TamingTF}, encoder/decoder models have been a core part of latent generative architectures, and serve to compress raw media (images, video, audio) into a lower-dimensional latent space. Latent diffusion models~\citep{rombach2021highresolution} typically use either a normal variational autoencoder (VAE)~\citep{kingma2013auto}, a quantized VAE such as a VQVAE~\citep{van2017neural} or VQGAN~\citep{Esser2021TamingTF} and its variants~\citep{lee2022autoregressive}, which adds a GAN discriminator loss~\citep{goodfellow_GAN} to achieve improved reconstruction quality with greater compression.  Most image and video generation models use convolutional autoencoders, though transformer-based encoding models such as Efficient-VQGAN~\citep{cao2023efficient}, ViT-VQGAN~\citep{yu2021vector}, and TiTok~\citep{yu2024image} show promising results using vision transformers. For the \OURS autoencoder, we found best results using a continuous convolutional VAE with discriminator loss.

Of the methods using convolutional autoencoders, the models have been split between image (2D) and video (3D) models. First are those which use an image VAE (with or without quantization) and encode frame-by-frame — for example Stable Video Diffusion~\citep{sdvideo}, Latent Shift~\citep{an2023latentshift}, VideoLDM~\citep{Blattmann2023AlignYL}, Emu-Video~\citep{emuvideo2023}, and CogVideo~\citep{hong2022cogvideo}. Models which use image encoders are typically unable to directly generate long, high FPS videos due to lack of temporal compression. A more recent alternative has been to use 3D or mixed 2D-3D models. For example, MAGViT~\citep{52431} uses a 3D VQGAN with both 3D and 2D downsampling layers, with average pooling for downsampling, and W.A.L.T.~\citep{gupta2023photorealistic} and MAGViT-V2~\citep{yu2023language} use fully 3D convolutional encoders with strided downsampling. In our work, we chose to use an interleaved 2D-1D (\eg, 2.5D) convolutional encoder, where we trade off a slight improvement to reconstruction quality from a fully 3D model for lower memory and computational costs.

A feature of W.A.L.T., MAGViT-V2 and also some ViT-based encoders such as C-ViViT~\citep{villegas2023phenaki} is the inclusion of causality, which is typically implemented for convolutions through an asymmetrical padding. Causality is usually implemented because the first frame is always encoded independently, allowing images to be explicitly encoded for joint image and video generation. However, we have found that causal encoding is not necessary to encode images and videos jointly — symmetrical padding functions for joint image and video generation, and encoded images with symmetrically-padded convolutions are able to be used as conditioning for image-to-video models. Symmetrical padding works for different video lengths as long as replicate padding is used; a similar result is reported by TATS~\citep{ge2022long}.
